% !TeX root = ../../brightPhD.tex
\chapter{总结与展望}
\label{chap:summary_outlook}

%\iffalse

\section{全文总结}
\label{sec:full_text_summary}

本文围绕连续体结构拓扑优化中的若干关键数值问题，针对传统低阶有限元离散方法在高精度分析、复杂局部约束处理以及高效计算实现方面所面临的局限，系统开展了相关数值方法设计、实现及其应用研究。全文按照“拓扑优化模型与数值基础--高精度离散方法研究--高性能软件平台实现”的总体思路展开，在变密度拓扑优化统一框架下，围绕任意次拉格朗日有限元、基于高阶有限元的多分辨率拓扑优化方法以及任意次胡张混合有限元等高精度数值方法，针对最小柔顺度、柔顺机构设计、近不可压缩材料与应力约束等典型问题，系统开展了离散建模、数值分析、算法构造与算例验证研究，并进一步面向上述模型与算法的统一实现需求，设计实现了高性能拓扑优化软件平台 SOPTX。相关研究贯通数学模型、离散格式、灵敏度分析、软件实现与数值验证等关键环节，形成了一条较为完整的研究链条，为连续体结构拓扑优化中高精度有限元离散、复杂约束建模与工程可复现计算提供了系统的理论、方法与实现支撑。

在高阶有限元离散研究方面，本文首先围绕单纯形单元与张量积单元上的任意次拉格朗日有限元，系统建立了线弹性问题的高阶离散格式，以及单元密度与节点密度两类设计变量表征下的拓扑优化模型与过滤策略。通过二维、三维典型结构及柔顺机构设计等数值算例，比较分析了不同单元类型、不同离散阶次以及不同密度表征方式对优化构型、力学响应精度与计算效率的影响，较为全面地揭示了任意次拉格朗日有限元在变密度拓扑优化中的数值特点及适用性。在此基础上，针对传统单分辨率框架中高精度物理分析与高分辨率拓扑描述难以兼顾的矛盾，本文进一步构建了基于高阶有限元的多分辨率高精度拓扑优化方法。该方法通过解耦设计网格与分析网格，建立了设计变量到物理密度的映射机制，给出了相应的刚度矩阵数值积分方法与灵敏度分析格式，并进一步推广至应力约束拓扑优化问题。数值结果表明，该方法能够在控制计算规模的同时兼顾结构分析精度与拓扑表达分辨率，为复杂约束条件下的高精度拓扑优化计算提供了一种有效途径。

针对近不可压缩材料与复杂局部约束条件下的拓扑优化建模问题，本文进一步研究了基于任意次胡张混合有限元的高精度拓扑优化方法。基于 Hellinger-Reissner 混合变分原理，系统建立了线弹性问题的应力--位移混合描述、任意次胡张混合有限元空间及其离散变分格式，并通过低阶稳定化处理与复杂边界条件下顶点应力连续性的部分松弛，提升了该方法在低阶情形与复杂工程边界下的数值适用性。在此基础上，围绕最小柔顺度、近不可压缩材料与局部应力约束等典型问题，分别构建了相应的拓扑优化列式与数值求解框架，形成了基于互补应变能的灵敏度解析格式、面向近不可压缩情形的插值模型修正策略以及基于独立应力的增广拉格朗日约束处理框架。数值结果表明，任意次胡张混合有限元方法在缓解体积闭锁、提高局部应力场逼近质量以及增强复杂约束评估可靠性方面具有明显优势，为近不可压缩材料与复杂局部约束条件下的高精度拓扑优化提供了新的建模与求解途径。

面向前述高精度离散模型与优化算法的统一实现需求，本文依托 FEALPy 开源智能 CAX 计算引擎，设计并实现了高性能拓扑优化软件平台 SOPTX。该平台继承 FEALPy 的张量抽象、多后端管理以及网格、函数空间、数值积分与线性代数求解等通用能力，在“分析--优化解耦”的设计原则下，围绕拓扑优化迭代闭环构建了面向领域问题的分层架构与核心模块体系，并实现了设计变量映射、结构响应分析、目标与约束评估、灵敏度分析及优化更新等关键环节的组件化封装与流程化组织。在关键技术层面，平台重点发展了基于解析预计算与张量收缩的高效矩阵组装与算子缓存机制、自动微分驱动的灵敏度分析实现，以及面向 NumPy、PyTorch、JAX 等不同后端与 CPU/GPU 异构设备的性能可移植执行机制。通过典型数值实验，进一步验证了 SOPTX 在高效矩阵组装、自动微分灵敏度分析以及多后端与异构设备切换方面的正确性、扩展性与计算性能，从而为本文提出的多类拓扑优化模型与高精度数值方法提供了统一、可复现且可扩展的软件实现载体。

总体来看，本文围绕连续体结构拓扑优化中的若干关键数值问题，在高精度离散建模、复杂局部约束处理与高效计算实现等方面开展了较为系统的研究，形成了贯通优化模型、离散格式、灵敏度分析与软件实现的研究体系。全文从典型拓扑优化问题出发，围绕高精度分析、复杂约束建模与统一软件实现三个层面，系统推进了相关理论、方法与实现工作之间的有机衔接。相关研究表明，将高精度有限元离散、复杂约束建模与高性能软件实现有机结合，是推动连续体结构拓扑优化向高精度、强约束与工程可复现方向发展的有效研究路径。


\section{研究展望}
\label{sec:research_outlook}

尽管本文围绕连续体结构拓扑优化中的若干关键数值问题，在高精度离散方法、多分辨率拓扑优化、胡张混合有限元建模与高性能软件平台实现等方面开展了较为系统的研究，但从连续体结构拓扑优化的发展趋势来看，相关工作仍有进一步深化和拓展的空间。总体而言，本文当前研究主要建立在线弹性、小变形与变密度拓扑优化框架之上，后续可围绕复杂物理与复杂约束下的高精度拓扑优化、高精度离散方法的进一步深化，以及高性能软件平台扩展与数据驱动智能计算方法融合三个方向继续开展研究。

首先，在复杂物理模型与复杂工程约束方面，本文现有的高阶有限元、多分辨率与胡张混合有限元方法仍有进一步推广空间。后续可在当前框架基础上，面向几何非线性、材料非线性、弹塑性与超弹性等问题，研究适用于大变形、复杂本构关系及路径依赖行为的高精度离散格式、灵敏度分析方法与稳健求解策略。同时，可进一步推广至热固耦合、流固耦合及热流固耦合等多物理场问题，以提升拓扑优化模型在复杂工况下的建模能力与分析精度。在约束处理方面，除本文已讨论的体积约束与应力约束外，后续还可研究屈曲约束、固有频率约束及增材制造相关约束，并结合大跨桥梁结构等更真实的工程设计场景，推动高精度拓扑优化方法由典型基准算例向复杂工程应用拓展。

其次，在高精度离散方法本身的深化方面，本文所研究的任意次拉格朗日有限元、基于高阶有限元的多分辨率高精度拓扑优化方法以及任意次胡张混合有限元方法，仍可在网格适应性、误差控制与计算效率等方面继续完善。后续可进一步推动多分辨率框架由当前较规则网格向非结构网格推广，并研究高阶有限元与混合有限元在更一般网格与复杂几何边界下的稳定离散与实现机制。同时，可将后验误差估计引入拓扑优化迭代过程，构建兼顾分析精度与计算代价的网格/阶次自适应策略，并结合移动网格方法提高结构边界、高应力区域及高梯度响应区域的局部分辨能力。对于胡张混合有限元方法，还可进一步探索自由度更低、计算代价更小而又保持应力逼近精度与稳定性的离散格式。此外，虚拟单元法作为另一类适用于一般多边形、多面体网格的高精度离散方法，后续也可考虑引入拓扑优化框架，以提升复杂几何与非结构网格条件下的离散灵活性和建模适应性。进一步地，本文现有的高精度离散、多分辨率设计与分析解耦以及复杂局部响应评估等思想，也可推广到相场方法、水平集方法等边界演化型拓扑优化框架中，从而推动高精度拓扑优化方法由变密度框架向更广义的拓扑优化方法拓展。

最后，在软件平台扩展与数据驱动智能计算方法融合方面，SOPTX 仍具有持续深化的空间。一方面，后续可围绕超大规模拓扑优化问题，研究大规模稀疏系统求解、高效并行组装、异构设备协同计算及分布式计算等关键技术，以提升平台在复杂工程问题中的可扩展性与计算效率。另一方面，可进一步完善其与 Abaqus、Ansys 等主流商业仿真软件之间的数据接口、模型转换与结果对接机制，建立统一的结果比对与交叉验证流程，从而增强平台的工程验证能力与实际应用价值。此外，可依托 FEALPy 与 SOPTX 已有的自动微分、多后端与统一实现框架，结合已有高精度数值算例数据，研究神经网络、深度学习等数据驱动智能计算方法在拓扑优化中的应用潜力与实现途径。通过上述拓展，SOPTX 有望逐步发展为兼具开放性、可扩展性、可复现性与一定数据驱动能力的高性能拓扑优化计算平台。


%\fi